\section{数据预处理}

\subsection{数据清洗}

首先对原始NSL-KDD数据集进行质量检查。经逐列扫描确认，41个特征列与1个标签列均不存在缺失值，因此无需进行填充或删除操作。但数据集中附带一列名为difficulty的字段，该字段为Tavallaee等人在构建NSL-KDD时引入的样本难度评分，与网络入侵检测任务本身无关，属于冗余信息，故予以删除。清洗后保留全部125,973条训练记录与22,544条测试记录，特征列从42列缩减为41列（其中38列为数值型，3列为分类型）。清洗后的数值特征整体分布情况如图~\ref{fig:ch2-numeric-dist}所示，三类分类特征的Top-10取值分布如图~\ref{fig:ch2-categorical}所示，可以观察到多数数值特征呈现明显的右偏长尾分布。

\subsection{分类特征编码}

NSL-KDD数据集包含三个字符串类型的分类特征：protocol\_type共3种取值（tcp、udp、icmp），其中tcp占据绝大多数；service共70余种取值，http最为常见；flag共11种取值，SF标志位占比最高。这种长尾分布的分类特征无法直接输入到数值型模型中，需要进行编码转换。

本文采用OneHot编码（OneHotEncoder）策略，将每个分类特征的每个取值转换为独立的二元特征列。具体而言，protocol\_type展开为3列，service展开为70列，flag展开为11列，去除原3个分类列后新增84列，最终特征维度从41维扩展至53维。这种编码方式完整保留了所有类别信息，避免了人为引入序数关系，对树模型与神经网络均友好。

\subsection{异常值处理}

网络流量数据中广泛存在离群点，而入侵检测任务中的离群点往往正是攻击流量的体现。若简单删除异常值所在行，会导致大量稀有攻击样本被错误清除，对检测性能产生负面影响。因此本文采用检测但不删除的策略：使用IQR方法（1.5倍四分位距）识别超出上下边界的取值，再使用clip截断策略将越界值替换为对应的边界值。

受异常值影响最严重的两个特征分别是srv\_diff\_host\_rate（异常率22.5\%）和dst\_host\_same\_src\_port\_rate（异常率19.9\%）。这两个特征本质上反映了扫描类攻击的统计指标，分布本身就具有显著的偏态特征，截断后仍能保留绝大部分判别信息。所有数值特征的相关系数热力图如图~\ref{fig:ch2-correlation}所示，可以观察到内容特征组内部存在较高相关性，而主机流量特征组与时间流量特征组之间也存在若干强相关对。

\subsection{长尾特征变换}

针对src\_bytes、dst\_bytes、count、srv\_count四个流量类特征的严重右偏分布，本文应用log1p变换（即$\log(1+x)$）进行压缩。这四个特征的原始取值跨度可达数个数量级，少数极大值会主导模型的梯度更新方向。log1p变换能够将偏态分布拉近正态，同时对零值保持兼容。变换前后的分布对比见图~\ref{fig:ch2-log1p}，原本集中在零附近且具有长尾的分布在变换后扩展为更接近正态的钟形分布。

\subsection{特征选择}

特征选择分两步执行，以在保留判别信息的同时降低维度、缓解过拟合。

第一步为方差阈值过滤。使用VarianceThreshold方法剔除方差接近零的特征列，主要清除OneHot编码后取值几乎全为零的稀疏列。该步骤将特征数从53维降至32维。

第二步为基于随机森林的特征重要度筛选。在训练集上拟合随机森林模型，按特征重要度从高到低排序，保留Top-20特征。最终特征重要度分布如图~\ref{fig:ch2-importance}所示。Top-5特征分别为：flag\_SF（重要度0.148）、same\_srv\_rate（0.108）、dst\_host\_srv\_count（0.105）、dst\_host\_same\_srv\_rate（0.081）、diff\_srv\_rate（0.076）。从排序可以看出，连接状态标志位以及目标主机的服务统计特征对入侵检测最具判别力，这与已有文献的结论一致。

\subsection{数据划分与持久化}

数据划分遵循NSL-KDD官方方案，训练集使用KDDTrain+.txt，测试集使用KDDTest+.txt，不进行随机重划分，以保证与已有研究结果的可比性。最终训练集形状为$(125973, 20)$，测试集形状为$(22544, 20)$，与原始数据集保持完全一致的样本规模。预处理过程中各阶段的特征维度变化路径如表~\ref{tab:dim-path}所示。

\begin{table}[htbp]
\centering
\caption{预处理各阶段特征维度变化}
\label{tab:dim-path}
\begin{tabular}{lcl}
\toprule
处理阶段 & 维度数 & 说明 \\
\midrule
原始数据 & 41 & 38数值特征 + 3分类特征 \\
OneHot编码 & 53 & 分类特征独热扩展 \\
异常值clip & 53 & 仅截断，维度不变 \\
log1p变换 & 53 & 对4个流量特征应用 \\
方差阈值过滤 & 32 & 移除低方差列 \\
RF Top-20筛选 & 20 & 按重要度排序取前20 \\
\bottomrule
\end{tabular}
\end{table}

处理完成后的数据以pickle格式持久化，保存于\texttt{outputs/processed/}目录下，共生成8个文件，分别对应训练集与测试集的特征矩阵与标签向量的二分类与五分类版本，便于后续模型加载使用。

\begin{figure*}[htbp]
\centering
\includegraphics[width=\textwidth]{04_numeric_distributions.png}
\caption{数值特征分布}
\label{fig:ch2-numeric-dist}
\end{figure*}

\begin{figure}[htbp]
\centering
\includegraphics[width=\columnwidth]{05_categorical_top10.png}
\caption{分类特征Top-10分布}
\label{fig:ch2-categorical}
\end{figure}

\begin{figure*}[htbp]
\centering
\includegraphics[width=\textwidth]{06_correlation_heatmap.png}
\caption{特征相关系数热力图}
\label{fig:ch2-correlation}
\end{figure*}

\begin{figure}[htbp]
\centering
\includegraphics[width=\columnwidth]{07_log1p_before_after.png}
\caption{log1p变换前后分布对比}
\label{fig:ch2-log1p}
\end{figure}

\begin{figure*}[htbp]
\centering
\includegraphics[width=\textwidth]{08_feature_importance_top20.png}
\caption{随机森林Top-20特征重要度}
\label{fig:ch2-importance}
\end{figure*}

经过上述预处理流程后，原始41维特征被压缩至20维关键特征，数据质量与可分性均得到明显改善，为下一章的模型构建奠定可靠的数据基础。第3章将基于本章产出的训练集与测试集，系统构建多种机器学习模型并进行对比实验。